\section{Typical codes on physical bit positions}\label{sec:native}

Run the simulator $N_s$ of \cref{thm:simulator} on a nonempty set
$G\subseteq G_{d,b}$. We will count inputs in $G$ with equal weight. The
construction has to keep track of which physical bit positions survive,
because the most frequently read bits can differ between traces.

\subsection{Counting the cells of a frozen trace}\label{sec:records}

For each attained full trace $\tau$, choose one set of revealed steps
$\Gamma_\tau$. Root the cone for input coordinate $q$ at the step that
appends output bit $\tpos(q)$. Delete both ports of each coordinate
whose input port has excessive degree. In the induced graph delete both
ports of every remaining coordinate whose paired distance is at most the
chosen radius. Denote the retained input positions by $U_\tau$ and set
$m_\tau=|U_\tau|$. The choices of steps and positions are common to the
whole trace. The existence of such cuts with bounded depth is proved in
\lean{TapeStorageRates.native_joint_typical_partition} using the walk theorem
of \cite[Corollary~O(a)]{PaperFourteen}.

For a fixed $\tau$, record the values at $\Gamma_\tau$ and the input bits
outside $U_\tau$. Call this complete record $z$. Let $G_{\tau,z}$ be
the inputs with that trace and record, and $C_{\tau,z}$ their
projections onto $U_\tau$. Fixing the exterior makes the projection
injective. It also turns every omitted helper into a constant, so
$C_{\tau,z}$ is a code for the induced relation on the retained ports.
Each codeword extends to a member of $G_{\tau,z}$.
These are the storage and extension statements of
\lean{TapeStorageRates.NativeTypicalCut.storage} and
\lean{TapeStorageRates.NativeTypicalCut.mem_code_iff}. The full trace fixes the successive append flags, hence the step at
which each output position is written.

\begin{proposition}[Counting native fibres]\label{prop:moment}
\formal{TapeStorageRates/NativeTypical.lean}{TapeStorageRates.NativeTypicalPartition.deficit_moment,TapeStorageRates.NativeTypicalPartition.tail_of_joint_count}
Let $X$ be uniform on $G$ and use a native partition as above. For an
attained cell define its deficit
$D_{\tau,z}=m_\tau-\log_2|C_{\tau,z}|$. Write $D_X$ for the deficit of
the cell containing $X$. For any real budget $\ell_{\tau,z}$ on each
attained record, with $\ell_X$ the corresponding random budget,
\[
 \mathbb E\,2^{D_X-\ell_X}
   =\frac{1}{|G|}\sum_\tau\sum_{z:\,G_{\tau,z}\ne\varnothing}
                2^{m_\tau-\ell_{\tau,z}} .
\]
If at most $2^h$ distinct pairs of trace and revealed step values occur
on $G$, then for every $t\ge0$,
\[
 \Pr\!\left[D_X>L-\log_2|G|+h+t\right]\le 2^{-t}.
\]
\end{proposition}
\begin{outline}
The $|G_{\tau,z}|$ inputs in a cell all contribute the same power of two.
Since $|C_{\tau,z}|=|G_{\tau,z}|$, this cardinality cancels when we sum
over the cell, giving its term in the identity.

For the second display, use the constant budget
$\ell=L-\log_2|G|+h+t$. A fixed trace and fixed revealed word allow at
most $2^{L-m_\tau}$ exterior assignments. Their combined contribution
to the unnormalized sum is at most $2^{L-\ell}$, independently of
$m_\tau$. Summing over the attained pairs gives
$\mathbb E\,2^{D_X-\ell}\le2^{-t}$. On the event $D_X>\ell$ its
integrand exceeds one. The same argument with the individual budgets
bounds $\Pr[D_X>\ell_X]$ when the set of possible exterior assignments
is known more precisely.
\end{outline}

\subsection{A tail bound at fixed degree}\label{sec:native-rate}

\begin{theorem}[Theorem N, typical native rate]\label{thm:native-tail}
\formal{TapeStorageRates/NativeTypicalRates.lean}{TapeStorageRates.typical_native_rate_tail,TapeStorageRates.typical_native_rates_full}
Fix $M$ with $k$ work tapes and $c\ge0$. Put $p=4k^2+1$ and
$Z_\Delta=2^{\sqrt{\log_2\Delta}}$. There are constants $K,A>0$ such
that at every sufficiently large integer cap $\Delta$ one can fix
an integer $s\ge1$ with
\[
       f_p(\Delta)\le 1/s\le K f_p(\Delta).
\]
For every fixed integer radius $r\ge0$, all sufficiently large $d$,
every $b\ge1$ and every nonempty $G\subseteq G_{d,b}$, the bit simulator
$N_s$ has a partition of $G$ with the following properties.

Projecting an attained cell as in \cref{sec:records} gives a storage code
on the retained input positions of its trace. Every
input and recovery port has degree at most $\Delta$, every surviving paired distance
exceeds $r$, and
\[
               m_\tau\ge (1-Z_\Delta/\Delta)L .
\]
The simulator halts within $(11c+1)(L+1)$ steps on $G$. For $X$ uniform
on $G$, let $R_X=\log_2|C_{\tau,z}|/m_\tau$ for its cell.
On this same partition, for every $u\ge0$,
\[
 \Pr\!\left[R_X<\frac{\log_2|G|}{L}-A f_p(\Delta)-u\right]
      \le 2^{-uL}.
\]
\end{theorem}
\begin{outline}
At fixed $s$ the startup cost fits in $L+1$ steps for all
sufficiently large $d$, uniformly in $b\ge1$. The exact clock and trace
count in \cref{thm:simulator} give the displayed time bound and at most
$2^{aL/s}$ traces on $G$, for a constant $a=a(M,c)$. All inputs in $G$
have the same output length.

On each trace, group the physical steps into the simulator's batches.
A complete batch has $10s+1$ steps. Each head sweeps all cells in three
aligned blocks and otherwise stays in those blocks. Replace its physical
position by its block index. The simultaneous work-head trajectory admits
nine coarse slots per complete batch and one for the final drain. This
is still a trajectory of $k$ unit-step work heads, with $O_c(L/s)$
vertices. Batch boundaries are retained, including when the heads stay
in the same blocks.

Apply the full-path cut theorem to this coarse trajectory with cut
fraction $\eta_s=\Theta_{M,c}(1/s)$
\cite[Definition~7.1 and Corollary~9.3(a)]{PaperFourteen}. Mark the
whole batch containing each selected coarse vertex, and reveal every
physical step record in the marked batches. A last-visit dependency
between distinct batches lifts to a coarse last-visit chain confined to
those endpoint batches: an intermediate complete batch visiting the same
block would also visit the precise physical cell, contrary to the
last-visit property. The chronological spine lifts as well. Counting
path lengths in vertices, a coarse free-path bound $H_s$ gives a physical bound
$(10s+1)H_s+1$. The extra vertex allows a revealed root; the suffix of
the path must avoid every marked batch. The path transfer is
\lean{TapeStorageRates.BatchCoarsening.Coarsening.path_length_le}.

We bound the joint support of these revelations using one local seed
per marked batch. At entry to a marked batch, the full simulator
control and the initial contents of its three-block input and work
windows determine every step record in that batch. The control includes
the buffered values and pending output. The seed uses $O_M(s)$ bits, so
a set $J$ of marked batches has at most $2^{O_M(s)|J|}$ joint record values;
no independence is assumed. Since $|J|\le O_c(\eta_s L/s)$, choosing
the constant in $\eta_s$ makes this support at most $2^{L/s}$. We fix
the complete records themselves. The seeds only bound their number.
The joint support estimate is
\lean{TapeStorageRates.NativeBatchRecord.compiler_joint_record_support}.

The coarse cut parameter is now $x=4/\eta_s=O_{M,c}(s)$. For the
physical depth parameter $D_s$, put
$B_s=\sum_{j=0}^{D_s}(k+1)^j$. The factor $10s+1$ in the lifted depth
preserves the estimate
\[
 \log_2\log_2(2sB_s)=O_{M,c}(s^p\log_2(s+1)).
\]
The numerical argument is formalized in
\lean{TapeStorageRates.genie_scaled_cap_degree_bound_at}. Choose the
constant in the inversion so that the double logarithm is at most
$\tfrac12\log_2\log_2\Delta$. Then $\max(s,B_s)\le Z_\Delta$ at the
stated scale for $s$, by
\lean{TapeStorageRates.native_batch_size_for_cap}. This gives the same
function $f_p$ as the packet bound while retaining the physical bit
positions and actual append roots.

There are at most $B_sL$ incidences. Removing input coordinates of degree
greater than $\Delta$ therefore removes at most $B_sL/(\Delta+1)$
coordinates. In the induced graph, the short-pair estimate bounds the
second deletion by $E_r L/\log_2 d$ for a constant $E_r$ depending on the
fixed parameters and radius, uniformly over the traces and the widths
$b\ge1$ \cite[Proposition~4.11(a), accompanying Lean theorem]{PaperEleven}.
Since $B_s/(\Delta+1)<Z_\Delta/\Delta$, both deletions fit in the
claimed fraction at all sufficiently large $d$. Deleting coordinates
can only increase the remaining paired distances.

Put $\lambda=Z_\Delta/\Delta$. For large caps $Z_\Delta^2\le\Delta$,
so $\lambda\le1/s$. The chosen traces and genie words jointly cost at
most $(a+1)L/s$ bits. Apply \cref{prop:moment} with $t=uL$. Outside
an event of probability at most $2^{-uL}$ we obtain
\[
 \log_2 |C_{\tau,z}|
 \ge\log_2|G|-\bigl(\lambda+(a+1)/s+u\bigr)L .
\]
The code is nonempty, so its rate on $m_\tau\le L$ coordinates is at
least $\log_2|C_{\tau,z}|/L$. The loss
$\lambda+(a+1)/s$ is at most $(a+2)K f_p(\Delta)$. This proves the
tail with a single partition for every confidence level.
\end{outline}

For a full successful cube $\log_2|G|/L=1$. Taking $u=1/s$ gives rate
$1-O_{M,c}(1/s)$ except on input mass at most $2^{-L/s}$. For a smaller
success set its own rate enters with coefficient one. The retained
positions and recovery relation may vary from trace to trace. The code can also vary with the revealed word and the exterior values.
The native coding target has the same rate-loss scale as the packet
target, while keeping the physical positions and append times available
for a more specific upper bound. Establishing such an upper bound remains open.

\subsection{A sufficient native coding bound}\label{sec:native-test}

The native half-rate hypothesis concerns the full cone relation of a
specified machine at a specified depth. It asks for an upper bound
$2^{(1/2+\delta)L}$ on every binary storage code on each eligible cone
relation for all sufficiently large $d$, for every fixed $\delta>0$.
The step width $\operatorname{blockWidth}(N_s,1)$ is the number of bits
charged for revealing one complete physical record.

\begin{corollary}[The native frozen-trace criterion]\label{cor:native}
\formal{TapeStorageRates/TraceTransfer.lean}{TapeStorageRates.native_frozen_trace_exclusion,TapeStorageRates.native_tests_exclude_linear}
For each $M$ and $c>0$ one can choose a single aligned simulator $N_s$ and
time coefficient $c'>0$ with the following implication. Suppose the
native half-rate hypothesis of \cite[Corollary L(b)]{PaperFourteen}
holds at canonical widths for $N_s$ at scale one, time coefficient
$c'$, and depth parameter $\operatorname{genieDepth}
(k,\operatorname{blockWidth}(N_s,1),\max(c',1),1/8)$. Then $M$
succeeds on the full cube at only finitely many canonical dimensions.
If this hypothesis holds on the chosen simulator for each source machine
and coefficient, \textnormal{(T)} is false.
\end{corollary}
\begin{outline}
An eligible run halts by the specified time bound. The hypothesis
ranges over cuts with the prescribed
remaining depth and all choices of times before halt at the prescribed
output-counter positions. In particular it includes the actual append
steps. Choose $s$ so that the complete trace support costs at most $L/4$
bits at large successful dimensions, using
\lean{TapeStorageRates.aligned_low_trace_dimensions}, which supplies
$c'=11c+1$. A largest trace
contains at least $2^{3L/4}$ inputs. Revealing at most $L/8$ bits on
that trace leaves a common word with at least $2^{5L/8}$ inputs. These
inputs form a code on the untrimmed append-time cones by the same
predecessor induction as in \cref{sec:model}. This exceeds the hypothesized
bound for every $0<\delta<1/8$. The formal criterion applies the skeleton-entropy lower bound of
Corollary~L(b). Skeleton entropy is at most the logarithm of the full
trace support. The separate witness theorem
\lean{TapeStorageRates.native_five_eighths} constructs the native code
and \lean{TapeStorageRates.NativeFibreWitness.outputTime_eq_appendTime}
identifies its roots.
\end{outline}

The same checked criterion allows any fixed $0<\eta<1/2$. Choose the
simulator so that its full trace support costs at most $\eta L$ bits
at large successful dimensions. In the depth parameter, replace the
cut fraction $1/8$ by $(1/2-\eta)/2$. The common trace and genie fibre
then has rate at least $1-\eta-(1/2-\eta)/2>1/2$.
